% Options for packages loaded elsewhere
\PassOptionsToPackage{unicode}{hyperref}
\PassOptionsToPackage{hyphens}{url}
\PassOptionsToPackage{dvipsnames,svgnames,x11names}{xcolor}
\documentclass[
  10pt,
  fontset=fandol]{ctexart}
\usepackage{xcolor}
\usepackage[margin=16mm]{geometry}
\usepackage{amsmath,amssymb}
\setcounter{secnumdepth}{-\maxdimen} % remove section numbering
\usepackage{iftex}
\ifPDFTeX
  \usepackage[T1]{fontenc}
  \usepackage[utf8]{inputenc}
  \usepackage{textcomp} % provide euro and other symbols
\else % if luatex or xetex
  \usepackage{unicode-math} % this also loads fontspec
  \defaultfontfeatures{Scale=MatchLowercase}
  \defaultfontfeatures[\rmfamily]{Ligatures=TeX,Scale=1}
\fi
\usepackage{lmodern}
\ifPDFTeX\else
  % xetex/luatex font selection
\fi
% Use upquote if available, for straight quotes in verbatim environments
\IfFileExists{upquote.sty}{\usepackage{upquote}}{}
\IfFileExists{microtype.sty}{% use microtype if available
  \usepackage[]{microtype}
  \UseMicrotypeSet[protrusion]{basicmath} % disable protrusion for tt fonts
}{}
\makeatletter
\@ifundefined{KOMAClassName}{% if non-KOMA class
  \IfFileExists{parskip.sty}{%
    \usepackage{parskip}
  }{% else
    \setlength{\parindent}{0pt}
    \setlength{\parskip}{6pt plus 2pt minus 1pt}}
}{% if KOMA class
  \KOMAoptions{parskip=half}}
\makeatother
\usepackage{longtable,booktabs,array}
\newcounter{none} % for unnumbered tables
\usepackage{calc} % for calculating minipage widths
% Correct order of tables after \paragraph or \subparagraph
\usepackage{etoolbox}
\makeatletter
\patchcmd\longtable{\par}{\if@noskipsec\mbox{}\fi\par}{}{}
\makeatother
% Allow footnotes in longtable head/foot
\IfFileExists{footnotehyper.sty}{\usepackage{footnotehyper}}{\usepackage{footnote}}
\makesavenoteenv{longtable}
\setlength{\emergencystretch}{3em} % prevent overfull lines
\providecommand{\tightlist}{%
  \setlength{\itemsep}{0pt}\setlength{\parskip}{0pt}}
\usepackage{amsmath,amssymb,mathtools,needspace}
\xeCJKDeclareCharClass{CJK}{"2032,"0400->"04FF,"2160->"216B,"2460->"2473,"25A0,"25C7,"25CB,"25CF}
\IfFontExistsTF{Arial Unicode MS}{\xeCJKsetup{AutoFallBack=true}\setCJKfallbackfamilyfont{\CJKrmdefault}{Arial Unicode MS}}{}
\setcounter{secnumdepth}{0}
\setlength{\emergencystretch}{2em}
\allowdisplaybreaks
\usepackage{bookmark}
\IfFileExists{xurl.sty}{\usepackage{xurl}}{} % add URL line breaks if available
\urlstyle{same}
\hypersetup{
  pdftitle={Newton-Cotes 公式},
  colorlinks=true,
  linkcolor={Maroon},
  filecolor={Maroon},
  citecolor={Blue},
  urlcolor={Blue},
  pdfcreator={LaTeX via pandoc}}

\title{Newton-Cotes 公式}
\author{}
\date{}

\begin{document}
\maketitle

\subsection{4.2 Newton-Cotes 公式}\label{newton-cotes-ux516cux5f0f}

\subsubsection{4.2.1 Cotes 系数}\label{cotes-ux7cfbux6570}

设将积分区间 \([a,b]\) 划分为 \(n\) 等份，步长 \(h=\dfrac{b-a}{n}\)，选取等距节点 \(x_k=a+kh\) 构

\href{../../source-images/pdf-096.jpeg}{原页图像}

造出的插值型求积公式

\[
I_n=(b-a)\sum_{k=0}^n C_k^{(n)}f(x_k)
\tag{4.2.1}
\]

称为 Newton-Cotes 公式，其中 \(C_k^{(n)}\) 称为 Cotes 系数。按式（4.1.6），引进变换

\[
x=a+th,
\]

则有

\[
C_k^{(n)}=\frac{h}{b-a}\int_0^n\prod_{\substack{j=0\\j\ne k}}^n\frac{t-j}{k-j}\,\mathrm{d}t
=\frac{(-1)^{n-k}}{nk!(n-k)!}\int_0^n\prod_{\substack{j=0\\j\ne k}}^n(t-j)\,\mathrm{d}t.
\tag{4.2.2}
\]

由于是多项式的积分，Cotes 系数的计算不会遇到实质性的困难。当 \(n=1\) 时，

\[
C_0^{(1)}=C_1^{(1)}=\frac12,
\]

这时求积公式是我们所熟悉的梯形公式（4.1.1）。

当 \(n=2\) 时，按式（4.2.2），Cotes 系数为

\[
C_0^{(2)}=\frac14\int_0^2(t-1)(t-2)\,\mathrm{d}t=\frac16,
\]

\[
C_1^{(2)}=-\frac12\int_0^2t(t-2)\,\mathrm{d}t=\frac46,\qquad
C_2^{(2)}=\frac14\int_0^2t(t-1)\,\mathrm{d}t=\frac16.
\]

相应的求积公式是下列 Simpson 公式：

\[
S=\frac{b-a}{6}\left[f(a)+4f\left(\frac{a+b}{2}\right)+f(b)\right].
\tag{4.2.3}
\]

而 \(n=4\) 的 Newton-Cotes 公式则特别称为 Cotes 公式，其形式为

\[
C=\frac{b-a}{90}[7f(x_0)+32f(x_1)+12f(x_2)+32f(x_3)+7f(x_4)].
\tag{4.2.4}
\]

这里 \(x_k=a+kh\)，\(h=\dfrac{b-a}{4}\)。

表 4.1 列出 Cotes 系数表开头的一部分。

表 4.1

{\def\LTcaptype{none} % do not increment counter
\begin{longtable}[]{@{}
  >{\raggedright\arraybackslash}p{(\linewidth - 12\tabcolsep) * \real{0.1429}}
  >{\raggedright\arraybackslash}p{(\linewidth - 12\tabcolsep) * \real{0.1429}}
  >{\raggedright\arraybackslash}p{(\linewidth - 12\tabcolsep) * \real{0.1429}}
  >{\raggedright\arraybackslash}p{(\linewidth - 12\tabcolsep) * \real{0.1429}}
  >{\raggedright\arraybackslash}p{(\linewidth - 12\tabcolsep) * \real{0.1429}}
  >{\raggedright\arraybackslash}p{(\linewidth - 12\tabcolsep) * \real{0.1429}}
  >{\raggedright\arraybackslash}p{(\linewidth - 12\tabcolsep) * \real{0.1429}}@{}}
\toprule\noalign{}
\begin{minipage}[b]{\linewidth}\raggedright
\(n\)
\end{minipage} & \begin{minipage}[b]{\linewidth}\raggedright
\(C_0^{(n)}\)
\end{minipage} & \begin{minipage}[b]{\linewidth}\raggedright
\(C_1^{(n)}\)
\end{minipage} & \begin{minipage}[b]{\linewidth}\raggedright
\(C_2^{(n)}\)
\end{minipage} & \begin{minipage}[b]{\linewidth}\raggedright
\(C_3^{(n)}\)
\end{minipage} & \begin{minipage}[b]{\linewidth}\raggedright
\(C_4^{(n)}\)
\end{minipage} & \begin{minipage}[b]{\linewidth}\raggedright
\(C_5^{(n)}\)
\end{minipage} \\
\midrule\noalign{}
\endhead
\bottomrule\noalign{}
\endlastfoot
1 & \(\frac12\) & \(\frac12\) & & & & \\
2 & \(\frac16\) & \(\frac23\) & \(\frac16\) & & & \\
3 & \(\frac18\) & \(\frac38\) & \(\frac38\) & \(\frac18\) & & \\
4 & \(\frac7{90}\) & \(\frac{16}{45}\) & \(\frac2{15}\) & \(\frac{16}{45}\) & \(\frac7{90}\) & \\
5 & \(\frac{19}{288}\) & \(\frac{25}{96}\) & \(\frac{25}{144}\) & \(\frac{25}{144}\) & \(\frac{25}{96}\) & \(\frac{19}{288}\) \\
\end{longtable}
}

\href{../../source-images/pdf-097.jpeg}{原页图像}

表 4.1（续表）

{\def\LTcaptype{none} % do not increment counter
\begin{longtable}[]{@{}
  >{\raggedright\arraybackslash}p{(\linewidth - 18\tabcolsep) * \real{0.1000}}
  >{\raggedright\arraybackslash}p{(\linewidth - 18\tabcolsep) * \real{0.1000}}
  >{\raggedright\arraybackslash}p{(\linewidth - 18\tabcolsep) * \real{0.1000}}
  >{\raggedright\arraybackslash}p{(\linewidth - 18\tabcolsep) * \real{0.1000}}
  >{\raggedright\arraybackslash}p{(\linewidth - 18\tabcolsep) * \real{0.1000}}
  >{\raggedright\arraybackslash}p{(\linewidth - 18\tabcolsep) * \real{0.1000}}
  >{\raggedright\arraybackslash}p{(\linewidth - 18\tabcolsep) * \real{0.1000}}
  >{\raggedright\arraybackslash}p{(\linewidth - 18\tabcolsep) * \real{0.1000}}
  >{\raggedright\arraybackslash}p{(\linewidth - 18\tabcolsep) * \real{0.1000}}
  >{\raggedright\arraybackslash}p{(\linewidth - 18\tabcolsep) * \real{0.1000}}@{}}
\toprule\noalign{}
\begin{minipage}[b]{\linewidth}\raggedright
\(n\)
\end{minipage} & \begin{minipage}[b]{\linewidth}\raggedright
\(C_0^{(n)}\)
\end{minipage} & \begin{minipage}[b]{\linewidth}\raggedright
\(C_1^{(n)}\)
\end{minipage} & \begin{minipage}[b]{\linewidth}\raggedright
\(C_2^{(n)}\)
\end{minipage} & \begin{minipage}[b]{\linewidth}\raggedright
\(C_3^{(n)}\)
\end{minipage} & \begin{minipage}[b]{\linewidth}\raggedright
\(C_4^{(n)}\)
\end{minipage} & \begin{minipage}[b]{\linewidth}\raggedright
\(C_5^{(n)}\)
\end{minipage} & \begin{minipage}[b]{\linewidth}\raggedright
\(C_6^{(n)}\)
\end{minipage} & \begin{minipage}[b]{\linewidth}\raggedright
\(C_7^{(n)}\)
\end{minipage} & \begin{minipage}[b]{\linewidth}\raggedright
\(C_8^{(n)}\)
\end{minipage} \\
\midrule\noalign{}
\endhead
\bottomrule\noalign{}
\endlastfoot
6 & \(\frac{41}{840}\) & \(\frac9{35}\) & \(\frac9{280}\) & \(\frac{34}{105}\) & \(\frac9{280}\) & \(\frac9{35}\) & \(\frac{41}{840}\) & & \\
7 & \(\frac{751}{17280}\) & \(\frac{3577}{17280}\) & \(\frac{1323}{17280}\) & \(\frac{2989}{17280}\) & \(\frac{2989}{17280}\) & \(\frac{1323}{17280}\) & \(\frac{3577}{17280}\) & \(\frac{751}{17280}\) & \\
8 & \(\frac{989}{28350}\) & \(\frac{5888}{28350}\) & \(\frac{-928}{28350}\) & \(\frac{10496}{28350}\) & \(\frac{-4540}{28350}\) & \(\frac{10496}{28350}\) & \(\frac{-928}{28350}\) & \(\frac{5888}{28350}\) & \(\frac{989}{28350}\) \\
\end{longtable}
}

可以看到，当 \(n\geq8\) 时，Cotes 系数有正有负，这时稳定性得不到保证。因此，实际计算中不用高阶的 Newton-Cotes 公式。

\begin{quote}
校注（agent 补充，CH04A-E01）：上述``当 \(n\geq8\) 时''若理解为对所有这类整数 \(n\) 成立，则存在原书表述错误。保留原文。由式（4.2.2）直接积分，\(n=9\) 时系数为 \[
\left(\frac{2857}{89600},\frac{15741}{89600},\frac{27}{2240},\frac{1209}{5600},\frac{2889}{44800},\frac{2889}{44800},\frac{1209}{5600},\frac{27}{2240},\frac{15741}{89600},\frac{2857}{89600}\right),
\] 全部为正；\(n=8\) 出现负系数这一事实仍成立。上述反例为 agent 按原公式进行的符号积分核验，不属于原书正文。
\end{quote}

\subsubsection{4.2.2 偶阶求积公式的代数精度}\label{ux5076ux9636ux6c42ux79efux516cux5f0fux7684ux4ee3ux6570ux7cbeux5ea6}

作为插值型的求积公式，\(n\) 阶的 Newton-Cotes 公式至少具有 \(n\) 次的代数精度（见定理 4.1）。实际的代数精度能否进一步提高呢？

先看 Simpson 公式（4.2.3），它是二阶 Newton-Cotes 公式，因此至少具有二次代数精度。进一步用 \(f(x)=x^3\) 进行检验，按 Simpson 公式计算，得

\[
S=\frac{b-a}{6}\left[a^3+4\left(\frac{a+b}{2}\right)^3+b^3\right].
\]

另一方面直接求积得 \(I=\int_a^b x^3\,\mathrm{d}x=\dfrac{b^4-a^4}{4}\)。这时有 \(S=I\)，即 Simpson 公式对次数不超过三次的多项式均能准确成立。又容易验证它对于 \(f(x)=x^4\) 通常是不准确的，因此，Simpson 公式实际上具有三次代数精度。

一般地，可以证明下述论断。

\textbf{定理 4.2} 当阶 \(n\) 为偶数时，Newton-Cotes 公式（4.2.1）至少有 \(n+1\) 次代数精度。

\textbf{证明} 只要验证，当 \(n\) 为偶数时，Newton-Cotes 公式对 \(f(x)=x^{n+1}\) 的余项为零。

按余项公式（4.1.7），由于这里 \(f^{(n+1)}(x)=(n+1)!\)，从而有

\[
R[f]=\int_a^b\prod_{j=0}^n(x-x_j)\,\mathrm{d}x.
\]

引进变换 \(x=a+th\)，并注意到 \(x_j=a+jh\)，有

\[
R[f]=h^{n+2}\int_0^n\prod_{j=0}^n(t-j)\,\mathrm{d}t.
\]

若 \(n\) 为偶数，则 \(\dfrac n2\) 为整数，再令 \(t=u+\dfrac n2\)，进一步有

\[
R[f]=h^{n+2}\int_{-n/2}^{n/2}\prod_{j=0}^n\left(u+\frac n2-j\right)\,\mathrm{d}u,
\]

\href{../../source-images/pdf-098.jpeg}{原页图像}

据此可以断定 \(R[f]=0\)，因为被积函数

\[
H(u)=\prod_{j=0}^n\left(u+\frac n2-j\right)=\prod_{j=-n/2}^{n/2}(u-j)
\]

是个奇函数。证毕。

\subsubsection{4.2.3 几种低阶求积公式的余项}\label{ux51e0ux79cdux4f4eux9636ux6c42ux79efux516cux5f0fux7684ux4f59ux9879}

首先考察梯形公式，按余项公式（4.1.7），梯形公式（4.1.1）的余项为

\[
R_{\mathrm{T}}=I-T=\int_a^b\frac{f''(\xi)}2(x-a)(x-b)\,\mathrm{d}x.
\]

这里积分的核函数 \((x-a)(x-b)\) 在区间 \([a,b]\) 上保号（非正），应用积分中值定理，在 \((a,b)\) 内存在一点 \(\eta\)，使

\[
R_{\mathrm{T}}=\frac{f''(\eta)}2\int_a^b(x-a)(x-b)\,\mathrm{d}x
=-\frac{f''(\eta)}{12}(b-a)^3,\quad\eta\in(a,b).
\tag{4.2.5}
\]

再研究 Simpson 公式（4.2.3）的余项 \(R_{\mathrm{S}}=I-S\)。为此构造次数不大于三的多项式 \(H(x)\)，使之满足

\[
H(a)=f(a),\quad H(b)=f(b),\quad H(c)=f(c),\quad H'(c)=f'(c),
\tag{4.2.6}
\]

这里 \(c=\dfrac{a+b}{2}\)。由于 Simpson 公式具有三次代数精度，它对于这样构造出的三次式 \(H(x)\) 是准确的：

\[
\int_a^b H(x)\,\mathrm{d}x=\frac{b-a}{6}[H(a)+4H(c)+H(b)].
\]

而由插值条件（4.2.6）知，上式右端实际上等于按 Simpson 公式（4.2.3）求得的积分值 \(S\)，因此积分余项

\[
R_{\mathrm{S}}=I-S=\int_a^b[f(x)-H(x)]\,\mathrm{d}x.
\]

不难证明，对于满足条件（4.2.6）的多项式 \(H(x)\)，其插值余项

\[
f(x)-H(x)=\frac{f^{(4)}(\xi)}{4!}(x-a)(x-c)^2(x-b),
\]

故有

\[
R_{\mathrm{S}}=\int_a^b\frac{f^{(4)}(\xi)}{4!}(x-a)(x-c)^2(x-b)\,\mathrm{d}x.
\]

这时积分核 \((x-a)(x-c)^2(x-b)\) 在 \([a,b]\) 上保号（非正），用积分中值定理有

\[
R_{\mathrm{S}}=\frac{f^{(4)}(\eta)}{4!}\int_a^b(x-a)(x-c)^2(x-b)\,\mathrm{d}x
=-\frac{b-a}{180}\left(\frac{b-a}{2}\right)^4f^{(4)}(\eta).
\tag{4.2.7}
\]

关于 Cotes 公式（4.2.4）的积分余项，这里不再具体推导，仅列出结果如下：

\[
R_{\mathrm{C}}=I-C=-\frac{2(b-a)}{945}\left(\frac{b-a}{4}\right)^6f^{(6)}(\eta).
\tag{4.2.8}
\]

\href{../../source-images/pdf-099.jpeg}{原页图像}

\subsubsection{4.2.4 复化求积法及其收敛性}\label{ux590dux5316ux6c42ux79efux6cd5ux53caux5176ux6536ux655bux6027}

前已指出，在使用 Newton-Cotes 公式时，提高阶的途径并不总能取得满意的效果。为了改善求积的精度，通常采用复化求积法。

设将积分区间 \([a,b]\) 划分为 \(n\) 等份，步长 \(h=\dfrac{b-a}{n}\)，分点为 \(x_k=a+kh\)，\(k=0,1,\cdots,n\)。所谓复化求积法，就是先用低阶的 Newton-Cotes 公式求得每个子区间 \([x_k,x_{k+1}]\) 上的积分值 \(I_k\)，然后再求和，用 \(\displaystyle\sum_{k=0}^{n-1}I_k\) 作为所求积分 \(I\) 的近似值。

复化梯形公式的形式是

\[
T_n=\sum_{k=0}^{n-1}\frac h2[f(x_k)+f(x_{k+1})]
=\frac h2\left[f(a)+2\sum_{k=1}^{n-1}f(x_k)+f(b)\right],
\tag{4.2.9}
\]

依式（4.2.5），其积分余项

\[
I-T_n=\sum_{k=0}^{n-1}\left[-\frac{h^3}{12}f''(\eta_k)\right]
=-\frac{b-a}{12}h^2f''(\eta).
\tag{4.2.10}
\]

此外，记子区间 \([x_k,x_{k+1}]\) 的中点为 \(x_{k+\frac12}\)，则复化 Simpson 公式为

\[
\begin{aligned}
S_n&=\sum_{k=0}^{n-1}\frac h6[f(x_k)+4f(x_{k+\frac12})+f(x_{k+1})]\\
&=\frac h6\left[f(a)+4\sum_{k=0}^{n-1}f(x_{k+\frac12})+2\sum_{k=1}^{n-1}f(x_k)+f(b)\right].
\end{aligned}
\tag{4.2.11}
\]

如果将每个子区间 \([x_k,x_{k+1}]\) 划分为 4 等分，内分点依次记作 \(x_{k+\frac14}\)，\(x_{k+\frac12}\)，\(x_{k+\frac34}\)，则复化 Cotes 公式具有形式

\[
\begin{aligned}
C_n=\frac h{90}\biggl[&7f(a)+32\sum_{k=0}^{n-1}f(x_{k+\frac14})+12\sum_{k=0}^{n-1}f(x_{k+\frac12})\\
&+32\sum_{k=0}^{n-1}f(x_{k+\frac34})+14\sum_{k=1}^{n-1}f(x_k)+7f(b)\biggr].
\end{aligned}
\tag{4.2.12}
\]

依据式（4.2.7）与式（4.2.8），容易得到复化 Simpson 公式和复化 Cotes 公式的余项，它们分别是

\[
I-S_n=-\frac{b-a}{180}\left(\frac h2\right)^4f^{(4)}(\eta),\quad\eta\in[a,b],
\tag{4.2.13}
\]

\[
I-C_n=-\frac{2(b-a)}{945}\left(\frac h4\right)^6f^{(6)}(\eta),\quad\eta\in[a,b].
\tag{4.2.14}
\]

其他 Newton-Cotes 公式亦可用类似的手续加以复化。显然，复化的 Newton-Cotes 公式仍然是形如式（4.1.3）的机械求积公式。

\textbf{例 4.1} 对于函数 \(f(x)=\dfrac{\sin x}{x}\)，试利用表 4.2 计算积分 \(I=\displaystyle\int_0^1\frac{\sin x}{x}\,\mathrm{d}x\)。

\textbf{解} 将积分区间 \([0,1]\) 划分为 8 等份，应用复化梯形法求得 \(T_8=0.945\,690\,9\)；而

\href{../../source-images/pdf-100.jpeg}{原页图像}

如果将 \([0,1]\) 划分为 4 等份，应用复化 Simpson 法有 \(S_4=0.946\,083\,2\)。

比较上面两个结果 \(T_8\) 与 \(S_4\)，它们都需要提供 9 个点上的函数值，计算量基本相同\footnote{使用求积公式时，计算的工作量主要耗费在函数求值上。}，然而精度却差别很大，同积分的准确值\footnote{所谓积分的``准确值''，是指每一位数字都是有效数字的积分值。} \(I=0.946\,083\,1\) 比较，复化梯形法的结果 \(T_8=0.945\,690\,9\) 只有两位有效数字，而复化 Simpson 法的结果 \(S_4=0.946\,083\,2\) 却有六位有效数字。

\begin{quote}
校注（agent 补充，CH04A-E02）：原文``只有两位有效数字''已保留。按本书 1.3.3 的式（1.3.2），以本页所列数值比较，\(|I-T_8|=0.0003922<0.0005\)，而大于 \(0.00005\)；因此 \(T_8\) 应具有三位有效数字。此为数值判断疑误，与原文转录状态分开记录。
\end{quote}

容易证明，复化的梯形法、Simpson 法和 Cotes 法当步长 \(h\to0\) 时均收敛到所求的积分值 \(I\)。现在考察当 \(h\) 很小时误差的渐近性态。

表 4.2

{\def\LTcaptype{none} % do not increment counter
\begin{longtable}[]{@{}ll@{}}
\toprule\noalign{}
\(x\) & \(f(x)\) \\
\midrule\noalign{}
\endhead
\bottomrule\noalign{}
\endlastfoot
\(0\) & \(1\) \\
\(1/8\) & \(0.997\,397\,8\) \\
\(1/4\) & \(0.989\,615\,8\) \\
\(3/8\) & \(0.976\,726\,7\) \\
\(1/2\) & \(0.958\,851\,0\) \\
\(5/8\) & \(0.936\,155\,6\) \\
\(3/4\) & \(0.908\,851\,6\) \\
\(7/8\) & \(0.877\,192\,5\) \\
\(1\) & \(0.841\,470\,9\) \\
\end{longtable}
}

先研究梯形法，按余项公式（4.2.10），有

\[
\frac{I-T_n}{h^2}=-\frac1{12}\sum_{k=0}^{n-1}hf''(\eta_k)\to-\frac1{12}\int_a^b f''(x)\,\mathrm{d}x,
\]

从而当 \(h\to0\) 时有下列渐近关系式：

\[
\frac{I-T_n}{h^2}\to-\frac1{12}[f'(b)-f'(a)].
\]

类似地，对于复化的 Simpson 法和 Cotes 法分别有

\[
\frac{I-S_n}{h^4}\to-\frac1{180\times2^4}[f'''(b)-f'''(a)],
\]

\[
\frac{I-C_n}{h^6}\to-\frac2{945\times4^6}[f^{(5)}(b)-f^{(5)}(a)].
\]

\textbf{定义 4.2} 如果一种复化求积公式 \(I_n\) 当 \(h\to0\) 时成立渐近关系式

\[
\frac{I-I_n}{h^p}\to C\quad(C\ne0,\text{定数}),
\]

则称求积公式 \(I_n\) 是 \(p\) 阶收敛的。

在这种意义下，复化的梯形法、Simpson 法和 Cotes 法分别具有二阶、四阶和六阶收敛精度。而当 \(h\) 很小时，对于复化的梯形法、Simpson 法和 Cotes 法分别有下列误差估计式：

\[
I-T_n\approx-\frac{h^2}{12}[f'(b)-f'(a)],
\tag{4.2.15}
\]

\[
I-S_n\approx-\frac1{180}\left(\frac h2\right)^4[f'''(b)-f'''(a)],
\tag{4.2.16}
\]

\[
I-C_n\approx-\frac2{945}\left(\frac h4\right)^6[f^{(5)}(b)-f^{(5)}(a)].
\tag{4.2.17}
\]

由此可见，若将步长 \(h\) 减半（即等份数 \(n\) 加倍），则梯形法、Simpson 法与 Cotes 法的误差分别减至原有误差的 \(1/4\)、\(1/16\) 与 \(1/64\)。

\begin{center}\rule{0.5\linewidth}{0.5pt}\end{center}

\subsection{相关原页校注（agent 补充）}\label{ux76f8ux5173ux539fux9875ux6821ux6ce8agent-ux8865ux5145}

以下校注来自本节覆盖的原页，可能也涉及同页相邻小节；原文照录与校正说明分开保留。

\subsubsection{原书 PDF 页 95 的校注}\label{ux539fux4e66-pdf-ux9875-95-ux7684ux6821ux6ce8}

\href{../pages/pdf-095.md}{完整原页转录} · \href{../../source-images/pdf-095.jpeg}{原页图像}

\begin{quote}
校注（agent 补充）：式（4.1.7）原扫描将导数阶数印作 \(f^{n+1}(\xi)\)，未带括号；此处按原式保留。由其引用的插值余项定理及上下文，这里表示 \(n+1\) 阶导数。
\end{quote}

\subsection{本节覆盖原页的校注记录}\label{ux672cux8282ux8986ux76d6ux539fux9875ux7684ux6821ux6ce8ux8bb0ux5f55}

下列记录按原页关联，含同页相邻小节的校注；计算时同时读取上文校注和完整原页。

\begin{itemize}
\tightlist
\item
  \texttt{CH04A-E01}：原书 PDF 页 97
\item
  \texttt{CH04A-E02}：原书 PDF 页 100
\end{itemize}

\end{document}
